% TOP_PROBLEM: {"id": 4800017, "title": "Multiple ergodic averages — Problem 17", "category_id": 11, "category": {"id": 11, "name": "dynamical_systems", "display_name": "Dynamical Systems", "description": "Problems about long-term behavior of deterministic systems, Hamiltonian dynamics, and periodic orbits.", "slug": "dynamical-systems", "order_index": 11, "created_at": "2026-07-31T15:26:25.671Z"}, "status": "partially_solved"}
% FIRST_POSTED: null
% ENTRY_KIND: statement_only
% Imported from: batch15.tex; UnsolvedMath ID 4800017
\documentclass[11pt]{article}
\usepackage[margin=25mm]{geometry}
\usepackage{fontspec}
\setmainfont{Latin Modern Roman}
\usepackage{amsmath,amssymb,mathtools}
\usepackage{unicode-math}
\setmathfont{Latin Modern Math}
\usepackage{xurl,hyperref}
\hypersetup{colorlinks=true,urlcolor=blue}
\setcounter{secnumdepth}{0}
\setlength{\parindent}{0pt}
\setlength{\parskip}{5pt}
\setlength{\emergencystretch}{3em}
\newcommand{\sref}[2]{\hyperlink{source:#1:#2}{[S#2]}}
\newcommand{\cataloguescope}[1]{\paragraph{Recorded target.}#1}
\begin{document}
% UnsolvedMath ID 4800017; source catalogue code AMR-047-0017
\setcounter{section}{4800016}
\section{Multiple ergodic averages — Problem 17}\label{4800017}
{\small\sffamily UnsolvedMath ID 4800017 · Dynamical Systems}
\subsection{Definitions and mathematical statement}

\paragraph{Shared notation.}$\mathbb N=\{1,2,\ldots\}$, and $\mathbb Z,\mathbb Q,\mathbb R,\mathbb C$ denote the integers, rationals, reals and complex numbers. Any other conventions are specified locally.


A system is a probability space $(X,\mathcal X,\mu)$ with invertible measure-preserving transformations; multiple transformations are assumed to commute. Write $T^j f=f\circ T^j$, and $T^j A$ for the image of a measurable set under $T^j$. A system with one transformation is ergodic if every invariant measurable set has measure zero or one. The space $L^\infty(\mu)$ consists of essentially bounded complex measurable functions. Fix $\ell\ge1$ and $p_1,\ldots,p_\ell\in\mathbb Z[t]$ with $p_j(0)=0$. Rational independence means that no nonzero rational linear combination of the $p_j$ is the zero polynomial; since the constant terms vanish, this is equivalent to saying every nontrivial such combination is nonconstant.
\paragraph{Statement recorded in the supplied source.}
For every commuting system, every $A\in\mathcal X$ and every $\varepsilon>0$, is there $n\in\mathbb N$ such that \[\mu\bigl(A\cap T_1^{p_1(n)}A\cap\cdots\cap T_\ell^{p_\ell(n)}A\bigr)\ge\mu(A)^{\ell+1}-\varepsilon?\] No ergodicity or distinct-degree hypothesis is imposed. \sref{4800017}{2}
\subsection{Short English statement}
Do rationally independent zero-constant polynomials force the optimal-power recurrence lower bound in every commuting system?
\subsection{Sources}
\begin{description}
\item[\hypertarget{source:4800017:1}{S1}] ULAM.AI and the UnsolvedMath Contributors, \emph{UnsolvedMath: A Curated Collection of Open Mathematics Problems}, record 4800017 (AMR-047-0017). CC BY 4.0. \url{https://huggingface.co/datasets/ulamai/UnsolvedMath}.
\item[\hypertarget{source:4800017:2}{S2}] N. Frantzikinakis, Some open problems on multiple ergodic averages (2016), Sections 1.2 and5.2.2--5.2.3, full Problem 17, equation (18), p.27. \url{https://arxiv.org/pdf/1103.3808v3}.
\end{description}
\label{end:4800017}
\end{document}
